\documentclass[11pt,a4paper]{article}

\usepackage[utf8]{inputenc}
\usepackage[T1]{fontenc}
\usepackage[french]{babel}
\usepackage{geometry}
\usepackage{hyperref}
\usepackage{listings}
\usepackage{xcolor}
\usepackage{booktabs}
\usepackage{array}
\usepackage{longtable}
\usepackage{enumitem}
\usepackage{fancyhdr}
\usepackage{parskip}

\geometry{margin=2cm, top=2.2cm, bottom=2.2cm}

\definecolor{codegreen}{rgb}{0,0.6,0}
\definecolor{codegray}{rgb}{0.5,0.5,0.5}
\definecolor{backcolour}{rgb}{0.97,0.97,0.97}
\definecolor{keywordcolor}{rgb}{0.13,0.29,0.53}
\definecolor{actionbg}{rgb}{0.95,0.98,0.95}

\lstdefinestyle{code}{
    backgroundcolor=\color{backcolour},
    commentstyle=\color{codegreen}\itshape,
    keywordstyle=\color{keywordcolor}\bfseries,
    basicstyle=\ttfamily\scriptsize,
    breaklines=true,
    frame=single,
    rulecolor=\color{black!25},
    numbers=none,
    tabsize=2
}
\lstset{style=code}

\pagestyle{fancy}
\fancyhf{}
\rhead{Résumé d'intégration \texttt{sendfd} -- OpenFaaS/faasd}
\lhead{Stage M2 ACS -- ISAE-SUPAERO}
\cfoot{\thepage}

\hypersetup{colorlinks=true, linkcolor=blue, urlcolor=cyan}

% Boite d'action
\newcommand{\act}[3]{%
  \vspace{0.2em}
  \noindent\colorbox{actionbg}{\parbox{\dimexpr\linewidth-2\fboxsep\relax}{%
    \small\textbf{#1~\texttt{#2}} \hfill \textit{\small #3}\\[-0.3em]
  }}
  \vspace{0.2em}
}

\title{
  \vspace{-1cm}
  \textbf{Intégration \texttt{sendfd} dans OpenFaaS/faasd}\\[0.3em]
  \large Résumé des modifications -- plan de mise en œuvre
}
\author{Stage M2 ACS -- ISAE-SUPAERO}
\date{\today}

\begin{document}
\maketitle
\thispagestyle{fancy}

% ================================================================
\section*{Contexte et objectif}
% ================================================================

L'objectif est d'éliminer les sauts TCP proxy (Gateway $\to$ faasd $\to$ Watchdog) en
transmettant le socket brut du client directement au container de la fonction via
\texttt{SCM\_RIGHTS} (sendfd/UDS). Le \textbf{faasd provider} est la seule autorité pour
résoudre l'IP du container. Chaque container crée un socket UDS identifié par son IP
tout en \textbf{conservant son serveur TCP habituel} (préservation de l'écosystème).
Un socket de relais gère les connexions HTTP keepalive (requête destinée à un autre container).

Activation : variable d'environnement \texttt{HTTPMIGRATE\_ENABLE=1} (Gateway et faasd)
et \texttt{SENDFD\_ENABLE=1} (watchdog). Sans ces variables : comportement original inchangé.

\bigskip\hrule\bigskip

% ================================================================
\section*{1. faasd provider -- 3 fichiers}
% ================================================================

\subsection*{1a. Nouveau fichier : \texttt{faasd/pkg/provider/handlers/ip\_endpoint.go}}

\textbf{Pourquoi :} La Gateway ne connaît pas l'IP d'un container. faasd est l'autorité
(accès containerd + CNI). Un endpoint dédié évite de modifier le type \texttt{FunctionStatus}
(contrat API public OpenFaaS préservé).

\textbf{Contenu :} une fonction exportée \texttt{MakeFunctionIPHandler(client *containerd.Client) http.HandlerFunc}.

Logique :
\begin{enumerate}[noitemsep]
  \item Lire \texttt{\{name\}} depuis l'URL (\texttt{mux.Vars(r)["name"]}).
  \item Appeler \texttt{GetFunction(client, name, namespace)} (existant dans
    \texttt{function\_get.go}).
  \item \texttt{GetFunction} appelle \texttt{cninetwork.GetIPAddress(name, task.Pid())}
    $\to$ retourne l'IP CNI du container.
  \item Écrire \texttt{\{"ip": fn.IP\}} en JSON / HTTP 200. Si absent : HTTP 404.
\end{enumerate}

\act{CRÉER}{faasd/pkg/provider/handlers/ip\_endpoint.go}{nouveau fichier, package handlers}

\subsection*{1b. Modifier : \texttt{faasd/cmd/provider.go}}

\textbf{Où :} dans \texttt{runProviderE()}, avant \texttt{bootstrap.Serve()}, ligne $\approx$105.

\textbf{Ajouter :}
\begin{lstlisting}[language=Go]
if os.Getenv("HTTPMIGRATE_ENABLE") == "1" {
    http.HandleFunc("/system/function-ip/",
        handlers.MakeFunctionIPHandler(client))
}
\end{lstlisting}

\textbf{Pourquoi :} enregistrement additionnel, aucun handler existant modifié.

\act{MODIFIER}{faasd/cmd/provider.go}{ajouter 4 lignes avant bootstrap.Serve()}

\subsection*{1c. Modifier : \texttt{faasd/pkg/provider/handlers/deploy.go}}

\textbf{Où :} dans \texttt{deploy()}, après \texttt{mounts := getOSMounts()}, ligne $\approx$140.

\textbf{Ajouter (conditionnel HTTPMIGRATE\_ENABLE=1) :} un bind-mount
\texttt{/var/lib/faasd/tlsmigrate} (hôte) $\to$ \texttt{/run/tlsmigrate} (container).

\textbf{Pourquoi :} le répertoire bind-monté est l'espace partagé entre Gateway et containers.
Le watchdog crée son socket UDS dans \texttt{/run/tlsmigrate/<ip>.sock} (côté container) ;
la Gateway y accède via \texttt{/var/lib/faasd/tlsmigrate/<ip>.sock} (côté hôte).
Mounts existants inchangés.

\act{MODIFIER}{faasd/pkg/provider/handlers/deploy.go}{ajouter bloc bind-mount conditionnel}

\bigskip\hrule\bigskip

% ================================================================
\section*{2. Gateway -- 6 fichiers}
% ================================================================

\subsection*{2a. Nouveau fichier : \texttt{faas/gateway/httpmigrate/loop.go}}

\textbf{Pourquoi :} \texttt{http.Server.ListenAndServe()} consomme les en-têtes HTTP avant
d'appeler le handler. Le FD transmis via sendfd serait alors tronqué. La loop intelligente
utilise \textbf{MSG\_PEEK} (lecture sans consommation) pour identifier la route et construire
le \texttt{*http.Request} sans toucher le buffer noyau.

\textbf{Fonctions :}
\begin{itemize}[noitemsep]
  \item \texttt{RunLoop(ln net.Listener, fallback *ChanListener, mux http.Handler)} :
    boucle Accept + goroutine par connexion.
  \item \texttt{handleConn(...)} : MSG\_PEEK $\to$ identifier route $\to$ branche
    \texttt{/function/} ou branche système.
  \item \texttt{serveFunctionConn(...)} : \texttt{http.ReadRequest(peeked)} $\to$ injecter
    rawFD dans contexte $\to$ \texttt{mux.ServeHTTP(rw, req)}.
\end{itemize}

\textbf{Résultat :} toute la pipeline OpenFaaS (auth, CallID, scale-from-zero, forwarding\_proxy)
s'exécute normalement via \texttt{mux.ServeHTTP}.

\act{CRÉER}{faas/gateway/httpmigrate/loop.go}{nouveau fichier, package httpmigrate}

\subsection*{2b. Nouveau fichier : \texttt{faas/gateway/httpmigrate/chan\_listener.go}}

Implémente \texttt{net.Listener} via un channel \texttt{chan net.Conn}.
Permet d'envoyer les connexions non-fonction (\texttt{/system/...}, \texttt{/healthz})
vers un \texttt{http.Server} de secours (même routeur \texttt{r}, mêmes timeouts).

\act{CRÉER}{faas/gateway/httpmigrate/chan\_listener.go}{nouveau fichier, package httpmigrate}

\subsection*{2c. Nouveau fichier : \texttt{faas/gateway/httpmigrate/conn\_rw.go}}

Implémente \texttt{http.ResponseWriter} écrivant dans \texttt{net.Conn}.
Méthode \texttt{SetMigrated()} : après sendfd réussi, \texttt{Write()} et
\texttt{WriteHeader()} ne font rien (la fonction répond directement au client).
Sans \texttt{SetMigrated()} (erreurs 401, 503) : écriture réelle sur TCP.

\act{CRÉER}{faas/gateway/httpmigrate/conn\_rw.go}{nouveau fichier, package httpmigrate}

\subsection*{2d. Nouveau fichier : \texttt{faas/gateway/httpmigrate/sendfd.go}}

Contient trois fonctions :
\begin{itemize}[noitemsep]
  \item \texttt{ResolveContainerIP(providerURL, funcName string)} :
    \texttt{GET http://faasd:8081/system/function-ip/\{name\}} $\to$ IP.
  \item \texttt{BuildSocketPath(dir, ip, suffix string)} :
    retourne \texttt{dir/ip+suffix+".sock"}.
  \item \texttt{SendFDs(sockPath string, clientFD, pipeWriteFD int)} :
    dial Unix $\to$ \texttt{WriteMsgUnix(SCM\_RIGHTS\{clientFD, pipeWriteFD\})}.
\end{itemize}

\act{CRÉER}{faas/gateway/httpmigrate/sendfd.go}{nouveau fichier, package httpmigrate}

\subsection*{2e. Nouveau fichier : \texttt{faas/gateway/httpmigrate/relay\_server.go}}

\textbf{Pourquoi :} HTTP keepalive -- une fonction peut recevoir une requête destinée à un
autre container sur le même socket TCP. Elle retourne le FD via ce socket de relais.

\textbf{Fonctions :}
\begin{itemize}[noitemsep]
  \item \texttt{CreateRelaySocket(hostDir, ip, providerURL string) (*RelayServer, error)} :
    crée \texttt{hostDir/<ip>-relay.sock}, lance goroutine d'écoute.
  \item Goroutine \texttt{listenLoop()} : reçoit le FD retourné via SCM\_RIGHTS, MSG\_PEEK
    $\to$ nouveau chemin $\to$ \texttt{ResolveContainerIP} $\to$ \texttt{SendFDs} vers le
    bon container.
\end{itemize}

Appelé depuis \texttt{sendfd.go / MigrateRequest()} juste avant \texttt{SendFDs()}.
Le container accède à ce socket via \texttt{/run/tlsmigrate/<own-ip>-relay.sock}
(même bind-mount).

\act{CRÉER}{faas/gateway/httpmigrate/relay\_server.go}{nouveau fichier, package httpmigrate}

\subsection*{2f. Modifier : \texttt{faas/gateway/handlers/forwarding\_proxy.go}}

\textbf{Où :} dans la closure de \texttt{MakeForwardingProxyHandler()}, au niveau du bloc
\texttt{forwardRequest()} + \texttt{time.Since(start)}.

\textbf{Ajouter la fonction \texttt{migrateRequest(w, r, rawFD, providerURL)} :}
\begin{enumerate}[noitemsep]
  \item \texttt{ResolveContainerIP()} $\to$ IP.
  \item \texttt{os.Pipe()} $\to$ \texttt{[pipeRead, pipeWrite]}, \texttt{t\_start}.
  \item \texttt{CreateRelaySocket()} + \texttt{SendFDs()}.
  \item \texttt{SetMigrated()} sur le ResponseWriter.
  \item Goroutine : \texttt{pipeRead.Read()} $\to$ \texttt{notifier.Notify("completed", realDuration)}.
\end{enumerate}

\textbf{Branchement :} si rawFD dans contexte ET \texttt{HTTPMIGRATE\_ENABLE=1} :
appeler \texttt{migrateRequest()} ; sinon : code vanilla original inchangé.

\textbf{Pourquoi :} \texttt{notifier.Notify("started")} reste à sa place.
\texttt{notifier.Notify("completed")} est appelé avec la durée réelle (via pipe).
Prometheus reçoit les métriques correctes. AlertManager non affecté.

\act{MODIFIER}{faas/gateway/handlers/forwarding\_proxy.go}{ajouter migrateRequest() + branchement}

\subsection*{2g. Modifier : \texttt{faas/gateway/main.go}}

\textbf{Où :} bloc final de \texttt{main()}, lignes $\approx$262--271.

\textbf{Modification :} envelopper \texttt{s.ListenAndServe()} dans \texttt{else} ;
ajouter branche \texttt{HTTPMIGRATE\_ENABLE=1} :
\begin{lstlisting}[language=Go]
if os.Getenv("HTTPMIGRATE_ENABLE") == "1" {
    ln, _ := net.Listen("tcp", fmt.Sprintf(":%d", tcpPort))
    fb := httpmigrate.NewChanListener()
    go (&http.Server{Handler: r, ...}).Serve(fb)
    httpmigrate.RunLoop(ln, fb, r) // bloque
} else {
    log.Fatal(s.ListenAndServe()) // original
}
\end{lstlisting}

\act{MODIFIER}{faas/gateway/main.go}{remplacer le bloc ListenAndServe() final}

\bigskip\hrule\bigskip

% ================================================================
\section*{3. of-watchdog -- 3 fichiers}
% ================================================================

\subsection*{3a. Modifier : \texttt{of-watchdog/config/config.go}}

\textbf{Où :} struct \texttt{WatchdogConfig} (fin de struct) et fonction \texttt{New()}.

\textbf{Ajouter dans la struct :}
\begin{lstlisting}[language=Go]
SendFDEnable    bool   // SENDFD_ENABLE=1
SendFDSocketDir string // SENDFD_SOCKET_DIR (defaut: /run/tlsmigrate)
\end{lstlisting}

\textbf{Ajouter dans \texttt{New()} :}
\begin{lstlisting}[language=Go]
cfg.SendFDEnable    = getBoolEnv(env, "SENDFD_ENABLE", false)
cfg.SendFDSocketDir = getEnv(env, "SENDFD_SOCKET_DIR", "/run/tlsmigrate")
\end{lstlisting}

\textbf{Pourquoi :} valeurs par défaut sûres (false). Tests existants non affectés.

\act{MODIFIER}{of-watchdog/config/config.go}{2 ajouts : struct + New()}

\subsection*{3b. Nouveau fichier : \texttt{of-watchdog/pkg/sendfd\_server.go}}

\textbf{Pourquoi :} le watchdog doit écouter \textbf{simultanément} sur TCP :8080 (existant,
inchangé) et sur son socket UDS (nouveau). Les deux serveurs sont indépendants.

\textbf{Fonctions :}
\begin{itemize}[noitemsep]
  \item \texttt{getContainerIP(iface string) (string, error)} :
    \texttt{net.InterfaceByName("eth0")} $\to$ première adresse IPv4 (sans masque CIDR).
    Pas de variable d'env injectée : lecture directe de l'interface réseau du container.
  \item \texttt{StartSendFDServer(cfg, handler) error} :
    \texttt{getContainerIP()} $\to$ créer \texttt{/run/tlsmigrate/<ip>.sock}
    $\to$ \texttt{os.Chmod(0o777)} $\to$ boucle Accept + goroutine.
  \item \texttt{handleFDMessage(conn, handler)} :
    \texttt{ParseUnixRights()} $\to$ \texttt{[clientFD, pipeWriteFD]}
    $\to$ relay vers \texttt{/run/tlsmigrate/<ip>-fn.sock} (socket de la fonction).
\end{itemize}

\act{CRÉER}{of-watchdog/pkg/sendfd\_server.go}{nouveau fichier, package watchdog}

\subsection*{3c. Modifier : \texttt{of-watchdog/pkg/watchdog.go}}

\textbf{Où :} dans \texttt{Start(ctx)}, après \texttt{http.HandleFunc(...)} ($\approx$ ligne 72)
et \textbf{avant} \texttt{listenUntilShutdown()}.

\textbf{Ajouter :}
\begin{lstlisting}[language=Go]
// Mode dual : UDS additionnel (TCP :8080 inchange apres)
if w.config.SendFDEnable {
    go func() {
        if err := StartSendFDServer(w.config, requestHandler); err != nil {
            log.Printf("[sendfd] %v", err)
        }
    }()
}
\end{lstlisting}

\textbf{Pourquoi :} la goroutine UDS est non-bloquante. \texttt{listenUntilShutdown()} démarre
normalement juste après. Health checks, métriques Prometheus, endpoint \texttt{/\_/health}
restent sur TCP :8080 et sont inchangés.

\act{MODIFIER}{of-watchdog/pkg/watchdog.go}{ajouter goroutine conditionnelle avant listenUntilShutdown()}

\bigskip\hrule\bigskip

% ================================================================
\section*{4. Processus de la fonction -- 1 fichier}
% ================================================================

\subsection*{4a. Adapter : \texttt{prototype/.../timing\_fn\_worker.c}}

\textbf{Base :} \texttt{prototype/faasd\_tlsmigrate/function/timing-fn/timing\_fn\_worker.c}

\textbf{Modifications par rapport au prototype existant :}
\begin{enumerate}[noitemsep]
  \item Créer \texttt{/run/tlsmigrate/<own-ip>-fn.sock} au démarrage
    (IP lue via \texttt{getifaddrs()} sur \texttt{eth0}).
  \item Recevoir \textbf{2 FDs} via \texttt{recvmsg(SCM\_RIGHTS)} : \texttt{clientFD}
    et \texttt{pipeWriteFD} (au lieu de 1 FD + payload TLS).
  \item Lire la requête HTTP depuis \texttt{clientFD} (buffer noyau intact).
  \item Écrire la réponse HTTP dans \texttt{clientFD} (réponse directe au client).
  \item Écrire 8 bytes (timestamp uint64 ns, \texttt{CLOCK\_MONOTONIC}) dans
    \texttt{pipeWriteFD}.
  \item \textbf{Cas keepalive :} MSG\_PEEK sur \texttt{clientFD} après réponse. Si une
    deuxième requête est présente et son chemin diffère : connecter à
    \texttt{/run/tlsmigrate/<own-ip>-relay.sock}, envoyer \texttt{clientFD} via
    SCM\_RIGHTS. La Gateway redirige vers le bon container.
\end{enumerate}

\act{MODIFIER}{prototype/.../timing\_fn\_worker.c}{adapter le prototype C existant}

\bigskip\hrule\bigskip

% ================================================================
\section*{Récapitulatif des fichiers}
% ================================================================

\begin{longtable}{lll}
\toprule
\textbf{Fichier} & \textbf{Action} & \textbf{Condition} \\
\midrule
\endhead
\texttt{faasd/pkg/provider/handlers/ip\_endpoint.go} & CRÉER & \texttt{HTTPMIGRATE\_ENABLE=1} \\
\texttt{faasd/cmd/provider.go} & MODIFIER & \texttt{HTTPMIGRATE\_ENABLE=1} \\
\texttt{faasd/pkg/provider/handlers/deploy.go} & MODIFIER & \texttt{HTTPMIGRATE\_ENABLE=1} \\
\texttt{faas/gateway/httpmigrate/loop.go} & CRÉER & \texttt{HTTPMIGRATE\_ENABLE=1} \\
\texttt{faas/gateway/httpmigrate/chan\_listener.go} & CRÉER & -- \\
\texttt{faas/gateway/httpmigrate/conn\_rw.go} & CRÉER & -- \\
\texttt{faas/gateway/httpmigrate/sendfd.go} & CRÉER & -- \\
\texttt{faas/gateway/httpmigrate/relay\_server.go} & CRÉER & -- \\
\texttt{faas/gateway/handlers/forwarding\_proxy.go} & MODIFIER & \texttt{HTTPMIGRATE\_ENABLE=1} \\
\texttt{faas/gateway/main.go} & MODIFIER & \texttt{HTTPMIGRATE\_ENABLE=1} \\
\texttt{of-watchdog/config/config.go} & MODIFIER & -- \\
\texttt{of-watchdog/pkg/sendfd\_server.go} & CRÉER & \texttt{SENDFD\_ENABLE=1} \\
\texttt{of-watchdog/pkg/watchdog.go} & MODIFIER & \texttt{SENDFD\_ENABLE=1} \\
\texttt{prototype/.../timing\_fn\_worker.c} & MODIFIER & -- \\
\midrule
\textbf{Total} & 5 MODIFIER + 8 CRÉER & \\
\bottomrule
\end{longtable}

\end{document}
